\documentclass{article}
\usepackage[utf8]{inputenc}
\usepackage{hyperref}
\usepackage{graphicx}
\usepackage{listings}
\usepackage{xcolor}

\title{Multi-GPU Acceleration of Homomorphic Encryption in FIDESlib}
\author{Halil Ibrahim Kanpak}
\date{January 2026}

\begin{document}

\maketitle

\section{Overview}
This project implements true multi-GPU parallelism for the FIDESlib Homomorphic Encryption (HE) library. By resolving bottlenecks related to static global state and device-specific caches, we enabled independent GPU contexts to operate simultaneously, achieving significant performance gains in HE matrix-vector multiplication.

\section{Hardware and Infrastructure}
The implementation was developed and tested on \textbf{RunPod} using an instance equipped with \textbf{2x NVIDIA L4 GPUs} (Ada Lovelace architecture, 24GB VRAM each). 

\subsection{Environment Setup}
\begin{itemize}
    \item \textbf{OS:} Ubuntu 22.04 LTS
    \item \textbf{CUDA:} 12.4
    \item \textbf{Toolkit:} OpenFHE (CKKS scheme) and FIDESlib
    \item \textbf{Communication:} NCCL 2.20
\end{itemize}

\section{Implementation Challenges}
The original FIDESlib architecture had several limitations that prevented scaling beyond a single GPU.

\subsection{Context Interference}
FIDESlib originally stored evaluation and rotation keys as static global variables. In a multi-GPU environment, this caused one GPU's keys to overwrite another's, leading to invalid memory accesses and kernel crashes.
\textbf{Solution:} Moved key storage to instance-level member variables within the \texttt{Context} class.

\subsection{CUDA Graph Cache Conflicts}
The library used CUDA graphs to accelerate repetitive NTT/INTT operations. However, the graph cache was only keyed by the size of the operation, ignoring the device. Since CUDA graphs are device-specific, executing a graph created on GPU 0 on GPU 1 resulted in failure.
\textbf{Solution:} Updated graph cache keys to a \texttt{std::pair<int, int>} containing both operations size and Device ID.

\subsection{Global Constants Initialization}
The NTT twiddle factors (\texttt{psi} tables) were only being copied to GPU 0 due to hardcoded device selection in \texttt{ConstantsGPU.cu}.
\textbf{Solution:} Implemented per-device initialization tracking and fixed device selection logic to ensure all GPUs receive necessary constants.

\section{Communication Layer: NCCL vs NVSHMEM}
While the initial proposal explored \textbf{NVSHMEM}, the final implementation adopted \textbf{NCCL} (NVIDIA Collective Communications Library). 
\textbf{Rationale:}
\begin{itemize}
    \item \textbf{Portability:} NCCL is more widely supported on cloud-provider environments like RunPod.
    \item \textbf{Reliability:} Encountered several context use problems and synchronization issues with NVSHMEM in multi-context scenarios.
    \item \textbf{Performance:} NCCL's \texttt{ncclAllReduce} provided efficient aggregation for the sums generated by parallel matrix-vector multiplication.
\end{itemize}

\section{Performance Results}
On 2x L4 GPUs, we achieved a \textbf{2.51x speedup} for a 8192 $\times$ 128 matrix multiplication.
\begin{itemize}
    \item \textbf{Single GPU:} 118.21 ms
    \item \textbf{Dual GPU:} 47.14 ms
\end{itemize}
The super-linear scaling is attributed to reduced memory pressure per GPU and better occupancy of individual GPU caches when handling smaller data partitions.

\section{Computation Graph for HE MatMul}
Below is the logical flow of the parallelized HE matrix-vector multiplication:

\begin{enumerate}
    \item \textbf{Host:} Generate CKKS Context and Keys.
    \item \textbf{Host:} Partition Matrix rows across $N$ GPUs.
    \item \textbf{Context Setup (Per GPU):} 
        \begin{itemize}
            \item Create independent FIDESlib Context.
            \item Load per-context Eval/Rotation Keys to GPU memory.
            \item Initialize device-specific NTT constants.
        \end{itemize}
    \item \textbf{Execution Loop:}
        \begin{itemize}
            \item \textbf{GPU 0..N:} Receive Ciphertext Matrix partition.
            \item \textbf{GPU 0..N:} Receive Ciphertext Vector (replicated).
            \item \textbf{GPU 0..N:} \texttt{multPt} (Hadamard product).
            \item \textbf{GPU 0..N:} \texttt{rescale}.
            \item \textbf{GPU 0..N:} \texttt{Rotate} + \texttt{Add} (Partial sum reduction).
        \end{itemize}
    \item \textbf{NCCL Layer:} \texttt{ncclAllReduce} to aggregate partial sums across all GPUs.
    \item \textbf{Host:} Decrypt and validate final result.
\end{enumerate}

\end{document}
